\section{Rotate a Linked List}
\label{sec:ll-rotate}

\problemstatement{Given the head of a linked list and an integer $k$, rotate
the list to the right by $k$ places and return the new head.}

\begin{patternbox}
Rotating right by $k$ is a \textbf{relink}, not a shift: make the list
circular, then break it at the right spot. Since rotating by the length
gives the same list, first reduce $k$ modulo the length.
\end{patternbox}

\subsection*{Circularise, then cut}

Walk to the tail, counting the length $n$, and connect the tail to the head
to form a ring. Reduce $k$ to $k \bmod n$. The new tail is the node at
position $n-k$ from the start; the node after it is the new head. Walk to
that new tail, set its \texttt{next} to null (breaking the ring), and return
the new head. This performs any rotation with a single relink, avoiding
$k$ separate moves.

\begin{ideabox}[Rotation Is One Cut on a Ring]
Joining the list into a circle makes every rotation just a choice of where
to break it. Reducing $k$ modulo $n$ handles $k\ge n$, and the break point
$n-k$ places the last $k$ nodes at the front.
\end{ideabox}

\subsection*{Algorithm}

\begin{enumerate}
  \item Find the length $n$ and the tail; connect tail to head (ring).
  \item $k \gets k \bmod n$; the new tail is $n-k$ steps from the head.
  \item Break after the new tail; return the node following it.
\end{enumerate}

\subsection*{Dry run}

\dryrun{$1\to2\to3\to4\to5$, $k=2$.}

\begin{itemize}
  \item $n=5$, $k=2$; new tail at position $5-2=3$ (node $3$).
  \item Break after $3$; new head $4$. Result $4\to5\to1\to2\to3$.
\end{itemize}

\subsection*{C++ implementation}

\begin{lstlisting}
#include <bits/stdc++.h>
using namespace std;

struct ListNode {
    int val; ListNode* next;
    ListNode(int v) : val(v), next(nullptr) {}
};

ListNode* build(const vector<int>& a) {
    ListNode dummy(0), *t = &dummy;
    for (int x : a) { t->next = new ListNode(x); t = t->next; }
    return dummy.next;
}

ListNode* rotateRight(ListNode* head, int k) {
    if (!head || !head->next || k == 0) return head;
    int n = 1;
    ListNode* tail = head;
    while (tail->next) { tail = tail->next; n++; }
    tail->next = head;                            // make it circular

    k %= n;
    int stepsToNewTail = n - k;
    ListNode* newTail = head;
    for (int i = 1; i < stepsToNewTail; i++) newTail = newTail->next;
    ListNode* newHead = newTail->next;
    newTail->next = nullptr;                       // break the ring
    return newHead;
}

void print(ListNode* h) { for (; h; h = h->next) cout << h->val << " "; cout << "\n"; }

int main() {
    print(rotateRight(build({1, 2, 3, 4, 5}), 2));   // 4 5 1 2 3
    return 0;
}
\end{lstlisting}

\begin{complexitybox}
\textbf{Time Complexity.} $O(n)$ -- one pass to measure, one to reach the
new tail. \textbf{Space Complexity.} $O(1)$.
\end{complexitybox}

\begin{takeawaybox}
Rotation is circularise-then-cut: one relink instead of $k$ moves, with
$k \bmod n$ collapsing large rotations. Turning a list into a ring to
simplify positional surgery is a handy general trick.
\end{takeawaybox}
